\reading{MR-5}{VaR Mapping}{grind}{4}
% ==========================================================

\begin{mrkeybox}[title={Learning objectives in this reading}]
\small
\begin{enumerate}[label=\textbf{\alph*.},leftmargin=1.7em]
\item Explain the principles underlying VaR mapping and describe the mapping process.
\item Explain how the mapping process captures general and specific risks, and calculate these
      risks in a portfolio given a set of primitive risk factors.
\item Differentiate among the three methods of mapping portfolios of fixed-income securities.
\item Summarize how to map a fixed-income portfolio into positions of standard instruments.
\item Describe how mapping of risk factors can support stress testing.
\item Explain how VaR can be calculated and used relative to a performance benchmark.
\item Describe the method of mapping forwards, forward rate agreements, interest rate swaps,
      and options.
\end{enumerate}
\end{mrkeybox}

% ---------------------------------------------------------
\lo{MR-5 a}{Explain the principles underlying VaR mapping and describe the mapping process.}

\begin{enumerate}[label=\textbf{\alph*)}]
\item Modern financial portfolios often include a vast array of assets such as bonds, stocks,
      currencies, commodities, and derivatives.
\item \term{Mapping} is the process of simplifying a large number of individual positions into a
      smaller number of aggregate positions based on elementary or primitive risk factors.
\end{enumerate}

\begin{mrexbox}[title={Examples}]
\begin{enumerate}
\item \term{Currency forwards.} A trader's desk with many dollar/euro forward contracts can
      aggregate these based on the major risk factor, the dollar/euro spot exchange rate.
\item \term{Bonds.} Map risks according to changing yields as bonds mature.
\item \term{Non-historical securities.} Map new assets like IPOs to existing risk factors.
\end{enumerate}
\end{mrexbox}

\begin{mrkeybox}[title={Benefits of VaR mapping}]
\begin{enumerate}[label=\textbf{\alph*)}]
\item \term{Simplification:} makes detailed data simpler.
\item \term{Efficiency:} focuses on fewer, significant risk factors to lessen computational load.
\item \term{Flexibility:} adapts to market or portfolio changes without needing extensive
      historical data.
\item \term{Predictive insight:} forecasts risks based on current and future market conditions.
\end{enumerate}
\end{mrkeybox}

% ---------------------------------------------------------
\lo{MR-5 b}{Explain how the mapping process captures general and specific risks, and calculate these risks in a portfolio given a set of primitive risk factors.}

\begin{center}
\small
\begin{tabularx}{\linewidth}{@{}p{4.6cm} >{\raggedright\arraybackslash}X >{\raggedright\arraybackslash}X@{}}
\arrayrulecolor{FRMRule}\toprule
\rowcolor{FRMNavy} \hdr{} & \hdr{What it is} & \hdr{Examples}\\
\midrule
\textbf{1. General risk factors}\newline (primitive risk factors)
  & Capture broad market movements that affect most assets in the portfolio to some degree.
  & Interest rates, equity market indices, etc.\\
\rowcolor{FRMSoft}
\textbf{2. Specific risk}\newline (unsystematic risk)
  & Arises from factors unique to individual assets or asset classes within the portfolio. It
    represents the variability \textit{not} explained by general risk factors.
  & Company-specific risks --- strikes, the PD of an individual bond, etc.\\
\bottomrule
\end{tabularx}
\end{center}

\begin{mrtrapbox}[title={The tradeoff}]
The mapping process aims to \term{minimize specific risk} by using a sufficient number of
general factors. The more factors you include, the more precisely you capture the overall risk
of the portfolio --- but the computation will be a lot. So keep a balanced set of general risk
factors.
\end{mrtrapbox}

% ---------------------------------------------------------
\lo{MR-5 c}{Differentiate among the three methods of mapping portfolios of fixed-income securities.}

\begin{center}
\small
\begin{tabularx}{\linewidth}{@{}p{1.9cm} >{\raggedright\arraybackslash}X >{\raggedright\arraybackslash}X >{\raggedright\arraybackslash}X@{}}
\arrayrulecolor{FRMRule}\toprule
\rowcolor{FRMNavy} \hdr{} & \hdr{Principal mapping} & \hdr{Duration mapping} & \hdr{Cash flow mapping}\\
\midrule
\textbf{What it does}
  & Focuses solely on the risk associated with the repayment of principal amounts.
  & Maps the risk of bonds to zero-coupon bonds with equivalent durations.
  & Decomposes bond risk into the risk associated with each bond's cash flows.\\
\rowcolor{FRMSoft}
\textbf{Accuracy}
  & Simplest approach among the three methods.
  & More accurate than principal mapping as it considers the timing of cash flows.
  & Most precise method but also the most complex, as it accounts for the timing and amount of
    cash flows.\\
\textbf{Catch}
  & \textit{Limitation:} ignores coupon payments, thus oversimplifying the risk assessment.
  & \textit{Challenge:} difficult to find a zero-coupon bond with an exact duration matching the
    portfolio.
  & \textit{Complexity:} high (incorporates correlations).\\
\bottomrule
\end{tabularx}
\end{center}

\begin{mrkeybox}[title={Remember the ordering}]
\begin{center}
\large $\mathrm{VaR}(\text{principal}) \;>\; \mathrm{VaR}(\text{duration}) \;>\;
\mathrm{VaR}(\text{diversified})$
\end{center}
\end{mrkeybox}

\begin{center}
\begin{tikzpicture}
\begin{axis}[
  width=10.4cm, height=4.6cm,
  xbar, bar width=16pt,
  xmin=2.50, xmax=3.10, ymin=-1.0, ymax=2.7,
  xlabel={\scriptsize VaR (\$ millions) --- note the axis starts at 2.50},
  xlabel style={font=\scriptsize}, tick label style={font=\scriptsize},
  xtick={2.5,2.6,2.7,2.8,2.9,3.0},
  ytick={0,1,2},
  yticklabels={{Cash flow\\(diversified)},{Duration},{Principal}},
  yticklabel style={align=right, font=\scriptsize},
  nodes near coords, nodes near coords style={font=\scriptsize\sffamily, color=black!75},
  every node near coord/.append style={/pgf/number format/fixed,
    /pgf/number format/precision=3},
  axis line style={draw=black!55}, clip=false]
\addplot[fill=PaleRed, draw=FRMRed!65, line width=0.4pt] coordinates {(2.968,2)};
\addplot[fill=PaleGold, draw=FRMGold!70, line width=0.4pt] coordinates {(2.737,1)};
\addplot[fill=PaleGreen, draw=FRMGreen!65, line width=0.4pt] coordinates {(2.615,0)};
\draw[black!35, {Stealth[length=3.5pt]}-{Stealth[length=3.5pt]}, line width=0.6pt]
  (axis cs:2.615,-0.34) -- (axis cs:2.968,-0.34);
\node[anchor=north, font=\tiny\sffamily, text=black!60, fill=white, inner sep=1.2pt]
  at (axis cs:2.79,-0.52) {a \$353k spread on identical positions};
\end{axis}
\end{tikzpicture}
\end{center}
\figcap{The same \$200 million portfolio, three mapping methods, three answers. The axis is
zoomed so the gaps are visible --- all three sit between \$2.6m and \$3.0m. Each step down
adds information: duration accounts for the \textit{timing} of cash flows and principal mapping
does not, and cash-flow mapping additionally accounts for the \textit{correlations} between
maturity buckets. More information means less apparent risk, which is why the ordering runs the
way it does --- and why it is a memorisable rule rather than a coincidence.}

% ---------------------------------------------------------
\lo{MR-5 d}{Summarize how to map a fixed-income portfolio into positions of standard instruments.}

\begin{mrexbox}[title={Example --- the same portfolio under all three methods}]
\textbf{Bond A:} 1 year, \$100 million, 3.5\% coupon. \quad
\textbf{Bond B:} 5 year, \$100 million, 5\% coupon.\\
Macaulay duration $= 2.768$. \quad Average maturity of bonds $= 3$ years.
\end{mrexbox}

\begin{center}
\small
\begin{tabular}{c r}
\arrayrulecolor{FRMRule}\toprule
\rowcolor{FRMNavy} \hdr{Maturity} & \hdr{VaR \%}\\
1 & 0.4696\\
\rowcolor{FRMSoft} 2 & 0.9868\\
3 & 1.4841\\
\rowcolor{FRMSoft} 4 & 1.9714\\
5 & 2.4261\\
\bottomrule
\end{tabular}
\end{center}

\begin{center}
\begin{tikzpicture}[every node/.style={font=\scriptsize\sffamily}]
  \draw[FRMNavy, {Stealth[length=5pt]}-{Stealth[length=5pt]}, line width=0.9pt]
    (0,0) -- (8.4,0);
  \draw[black!70, line width=0.8pt] (0.9,-0.16) -- (0.9,0.16);
  \draw[black!70, line width=0.8pt] (7.5,-0.16) -- (7.5,0.16);
  \draw[FRMGold, line width=1.1pt] (6.0,-0.22) -- (6.0,0.22);
  \node[anchor=north, text=black] at (0.9,-0.24) {2 year};
  \node[anchor=north, text=black] at (7.5,-0.24) {3 year};
  \node[anchor=north, text=FRMGold] at (6.0,-0.24) {\textbf{2.768 years}};
  \node[anchor=south, text=black] at (0.9,0.26) {VaR 0.9868};
  \node[anchor=south, text=black] at (7.5,0.26) {VaR 1.4841};
\end{tikzpicture}
\end{center}
\figcap{The 2.768-year duration falls between the 2-year and 3-year grid points, so its VaR\,\%
has to be interpolated between the two published figures rather than read off directly.}

The VaR of a 2.768-year maturity zero-coupon bond is interpolated as follows:
\[ 0.9868 + (1.4841 - 0.9868) \times (2.768 - 2) \;=\; 0.9868 + (0.4973 \times 0.768)
   \;=\; \mathbf{1.3687} \]

\subsection{1) Principal mapping calculation}
\begin{mrfmlbox}
\[ \mathrm{VaR}(\text{portfolio}) \;=\; NP \times \mathrm{VaR}(\text{average maturity of bonds}) \]
\mrvk{$NP$ = notional principal of the portfolio;\quad the VaR\,\% is read off at the
\textit{average maturity}, here 3 years.}{}
\end{mrfmlbox}
\[ \text{Principal mapping VaR} \;=\; \$200 \text{ million} \times 1.4841\% \;=\;
   \mathbf{\$2.968 \text{ million}} \]

\subsection{2) Duration mapping calculation}
\begin{mrfmlbox}
\[ \mathrm{VaR}(\text{portfolio}) \;=\; \mathrm{VaR}(\text{interpolated}) \times NP \]
\mrvk{the interpolated VaR\,\% is taken at the portfolio's \textit{duration}, here 2.768 years.}{}
\end{mrfmlbox}
\[ \text{Duration mapping VaR} \;=\; \$200 \text{ million} \times 1.3687\% \;=\;
   \mathbf{\$2.737 \text{ million}} \]

\subsection{3) Cash flow mapping}
\begin{mrfmlbox}
\[ \text{Undiversified VaR} \;=\; \sum \left( PV(\text{cash flow}) \times \mathrm{VaR}\,\% \right) \]
\mrvk{Diversified VaR additionally considers the correlations between the maturity buckets.}{}
\end{mrfmlbox}

\begin{center}
\scriptsize
\begin{tabular}{c r r r r r r r r}
\arrayrulecolor{FRMRule}\toprule
\rowcolor{FRMNavy} \hdr{Year} & \hdr{$x$} & \hdr{$x \times V$} & \hdr{1Y} & \hdr{2Y} & \hdr{3Y}
  & \hdr{4Y} & \hdr{5Y} & \hdr{$x\Delta$VaR}\\
\multicolumn{3}{c}{} & \multicolumn{5}{c}{\scriptsize\textit{Correlation matrix (R)}} & \\
1 & 104.83 & 0.4923 & 1 & 0.894 & 0.887 & 0.871 & 0.861 & 1.17\\
\rowcolor{FRMSoft} 2 & 4.63 & 0.0457 & 0.894 & 1 & 0.99 & 0.964 & 0.954 & 0.115\\
3 & 4.42 & 0.0656 & 0.887 & 0.99 & 1 & 0.992 & 0.987 & 0.170\\
\rowcolor{FRMSoft} 4 & 4.24 & 0.0836 & 0.871 & 0.964 & 0.992 & 1 & 0.996 & 0.217\\
5 & 81.88 & 1.9864 & 0.861 & 0.954 & 0.987 & 0.996 & 1 & 5.168\\
\midrule
\multicolumn{2}{r}{\textbf{Undiversified VaR}} & \textbf{2.674} & & & & & & 6.840\\
\multicolumn{2}{r}{\textbf{Diversified VaR}} & \textbf{2.615} & & & & & & \\
\bottomrule
\end{tabular}
\end{center}
\figcap{Each year's present-valued cash flow $x$ is multiplied by that maturity's VaR\,\% to give
$x \times V$. Summing that column ignores correlation and gives the undiversified VaR; applying
the correlation matrix instead gives the smaller diversified figure.}

% ---------------------------------------------------------
\lo{MR-5 e}{Describe how mapping of risk factors can support stress testing.}

\begin{enumerate}[label=\textbf{\alph*)}]
\item Stress testing a portfolio with the assumption of \term{perfect correlation} allows for a
      simpler evaluation of potential risks without considering the complexity of correlation
      effects.
\item However, this method is \term{not accurate} if actual correlations are not perfect.
\end{enumerate}

\begin{mrkeybox}[title={Just remember}]
When assuming perfect correlation (correlation $= 1$), the portfolio VaR is equal to the
\term{sum of the individual VaRs} --- the undiversified VaR.
\end{mrkeybox}

% ---------------------------------------------------------
\lo{MR-5 f}{Explain how VaR can be calculated and used relative to a performance benchmark.}

\subsection{Benchmarking a portfolio}
Benchmarking VaR compares the VaR of your portfolio to a reference portfolio, typically an
index. This helps assess the risk of your portfolio relative to a standard.

\textbf{Steps involved:}
\begin{enumerate}
\item \term{Match durations.} Ensure your portfolio and the benchmark portfolio have similar
      durations to compare risks from interest rate changes. This might involve adjusting
      weightings in your portfolio.
\item \term{Calculate absolute VaR.} Multiply the VaR percentages (the risk measure for
      individual holdings) by the market value of each asset in your portfolio. Sum these
      products to get the total VaR for your portfolio. Do the same for the benchmark.
\item \term{Calculate tracking error VaR.} This measures the difference between your portfolio's
      VaR and the benchmark's VaR. It reflects how closely your portfolio's VaR tracks the
      benchmark's VaR due to interest rate movements.
\end{enumerate}

Tracking error can be used to compute the variance reduction (similar to R-squared in a
regression) as follows:

\begin{mrfmlbox}
\[ \text{Variance improvement} \;=\; 1 - \left( \frac{\text{tracking error}}
   {\text{benchmark VaR}} \right)^{\!2} \]
\mrvk{tracking error = the VaR of the difference between the portfolio and the benchmark;\quad
benchmark VaR = the absolute VaR of the reference portfolio.}{}
\end{mrfmlbox}

\begin{mrkeybox}
The higher this ratio, the lower the tracking error VaR --- which is actually good.
\end{mrkeybox}

% ---------------------------------------------------------
\lo{MR-5 g}{Describe the method of mapping forwards, forward rate agreements, interest rate swaps, and options.}

\subsection{Mapping approaches for linear derivatives}

\subsubsection{1) Forwards and futures}
These are the simplest types of derivatives. Since their value is linear in the underlying spot
rates, their risk can be constructed easily from basic building blocks. We can use the
\term{delta-normal VaR} calculation here for linear portfolios.

\begin{mrexbox}[title={Example --- diversified VaR of a currency forward}]
Suppose you wish to compute the diversified VaR of a forward contract that is used to purchase
100 million euros in exchange for \$126.5 million one year from now. This forward position is
analogous to the following three separate risk positions:
\begin{enumerate}
\item A short position in a U.S.\ Treasury bill.
\item A long position in a one-year euro bill.
\item A long position in the euro spot market.
\end{enumerate}
\end{mrexbox}

\begin{center}
\begin{tikzpicture}[every node/.style={font=\scriptsize\sffamily},
  legbox/.style={rectangle, rounded corners=2.5pt, line width=0.8pt,
    text width=3.05cm, align=center, inner sep=6pt, minimum height=1.35cm}]
  \node[legbox, draw=FRMGold!70, fill=FRMPaper, text width=4.2cm, minimum height=1.1cm]
    (fwd) at (0,1.9) {\textbf{1-year forward}\\ buy EUR 100m for \$126.5m};
  \node[legbox, draw=FRMRed!60, fill=PaleRed] (l1) at (-4.0,0)
    {\textbf{(1) Short USD bill}\\ borrow dollars today};
  \node[legbox, draw=FRMNavy!60, fill=PaleNavy] (l2) at (0,0)
    {\textbf{(2) Long EUR bill}\\ lend euros for one year};
  \node[legbox, draw=FRMGreen!60, fill=PaleGreen] (l3) at (4.0,0)
    {\textbf{(3) Long EUR spot}\\ convert dollars to euros};
  \foreach \t in {l1,l2,l3}{
    \draw[FRMGold, -{Stealth[length=4.5pt]}, line width=0.9pt] (fwd) -- (\t);}
  \node[anchor=north, text=black!60, text width=11cm, align=center] at (0,-0.95)
    {each leg is then mapped to its own primitive risk factor};
\end{tikzpicture}
\end{center}
\figcap{The single forward contract is broken into the three primitive exposures it is
economically equivalent to, and each is then mapped to its own risk factor.}

\begin{center}
\scriptsize
\begin{tabular}{l r r r r r}
\arrayrulecolor{FRMRule}\toprule
\rowcolor{FRMNavy} \hdr{Risk factor} & \hdr{Price/Rate} & \hdr{VaR\%} & \hdr{EUR Spot}
  & \hdr{1Yr EUR} & \hdr{1Yr US}\\
\multicolumn{3}{c}{} & \multicolumn{3}{c}{\scriptsize\textit{Correlation matrix}}\\
EUR spot & 1.2500 & 4.5381 & 1.000 & 0.115 & 0.073\\
\rowcolor{FRMSoft} Long EUR bill & 0.0170 & 0.1396 & 0.115 & 1.000 & $-0.047$\\
Short USD bill & 0.0292 & 0.2121 & 0.073 & $-0.047$ & 1.000\\
\rowcolor{FRMSoft} EUR forward & 1.2650 & & & & \\
\bottomrule
\end{tabular}
\hspace{6pt}
\begin{tabular}{l r r r r r}
\arrayrulecolor{FRMRule}\toprule
\rowcolor{FRMNavy} \hdr{Position} & \hdr{PV factor} & \hdr{CF} & \hdr{$x$} & \hdr{$|x_i|\,V_i$}
  & \hdr{$x\Delta$VaR}\\
EUR spot & & & 122.911 & 5.578 & 31.116\\
\rowcolor{FRMSoft} Long EUR bill & 0.9777 & 100.0 & 122.911 & 0.172 & 0.142\\
Short USD bill & 0.9678 & 126.5 & 122.911 & 0.261 & $-0.036$\\
\midrule
\textbf{Undiversified VaR} & & & & \textbf{6.010} & 31.221\\
\bottomrule
\end{tabular}
\end{center}
\figcap{Left: the monthly VaR of each primitive risk factor and how the three factors correlate.
Right: the same three legs valued and summed. The $|x_i|V_i$ column adds to the undiversified
VaR.}

\begin{mrkeybox}[title={Note}]
The volatility of the \term{spot contract} is the highest by far. This is much greater than the
VaR for the EUR 1-year bill and the VaR for the USD bill.
\end{mrkeybox}

\subsubsection{2) Forward rate agreements}
\begin{enumerate}[label=\textbf{\alph*)}]
\item FRAs are forward contracts that allow users to lock in an interest rate at some future
      date.
\item The \term{buyer} of an FRA locks in a borrowing rate; the \term{seller} locks in a lending
      rate. In other words, the ``long'' receives a payment if the spot rate is above the
      forward rate.
\end{enumerate}

\begin{mrexbox}[title={Example --- VaR of an FRA}]
Suppose that you sold a $6 \times 12$ FRA on \$100 million. This is equivalent to borrowing
\$100 million for 6 months and investing the proceeds for 12 months. Assuming that the 360-day
spot rate is 4.5\% and the 180-day spot rate is 4.1\%, the present value of the notional
\$100 million contract is
\[ x \;=\; \frac{\$100}{1.0205} \;=\; \$97.991 \text{ million} \]
\end{mrexbox}

\begin{center}
\small
\begin{tabular}{@{}p{7.4cm}@{}}
\arrayrulecolor{FRMRule}\toprule
\rowcolor{PaleNavy}
\textbf{Short $6\times12$ FRA} = short/borrow 6-month bill \;+\; long/invest 12-month bill\\
\midrule
\rowcolor{PaleNavy}
\textbf{Long $6\times12$ FRA} = long/invest 6-month bill \;+\; short/borrow 12-month bill\\
\bottomrule
\end{tabular}
\end{center}

\begin{center}
\small
\begin{tabular}{l r r r r r r}
\arrayrulecolor{FRMRule}\toprule
\rowcolor{FRMNavy} \hdr{Position} & \hdr{PV(CF), $x$} & \hdr{VaR\%} & \hdr{} & \hdr{}
  & \hdr{$|x_i|\,V_i$} & \hdr{$x\Delta$VaR}\\
\multicolumn{3}{c}{} & \multicolumn{2}{c}{\scriptsize\textit{Correlations (R)}} & & \\
180 days & $-97.991$ & 0.1629 & 1 & 0.79 & 0.160 & $-0.0325$\\
\rowcolor{FRMSoft} 360 days & 97.991 & 0.4696 & 0.79 & 1 & 0.460 & 0.1537\\
\midrule
\textbf{Undiversified VaR} & & & & & \textbf{0.620} & 0.1212\\
\textbf{Diversified VaR} & & & & & \textbf{0.348} & \\
\bottomrule
\end{tabular}
\end{center}
\figcap{The two legs are of opposite sign and 0.79 correlated, so most of the risk cancels: the
diversified VaR of 0.348 is barely half the undiversified 0.620.}

\subsubsection{3) Interest rate swaps}
Interest rate swaps are commonly used to exchange interest rates from fixed to floating rates or
from floating to fixed rates. Thus, an interest rate swap can be broken down into fixed and
floating parts. The \term{fixed part is priced with a coupon-paying bond} and the \term{floating
part is priced as a floating-rate note}.

\begin{center}
\begin{tikzpicture}[every node/.style={font=\scriptsize\sffamily},
  legbox/.style={rectangle, rounded corners=2.5pt, line width=0.8pt,
    text width=4.6cm, align=center, inner sep=6pt, minimum height=1.25cm}]
  \node[legbox, draw=FRMGold!70, fill=FRMPaper, text width=4.0cm, minimum height=1.25cm]
    (sw) at (-4.4,0) {\textbf{4-year swap}\\ pay fixed, receive floating};
  \node[legbox, draw=FRMRed!60, fill=PaleRed] (fx) at (2.1,1.0)
    {\textbf{Short} a coupon-paying bond\\ \textit{(the fixed leg you pay)}};
  \node[legbox, draw=FRMGreen!60, fill=PaleGreen] (fl) at (2.1,-1.0)
    {\textbf{Long} a floating-rate note\\ \textit{(the floating leg you receive)}};
  \draw[FRMGold, -{Stealth[length=4.5pt]}, line width=0.9pt] (sw.east) -- (fx.west);
  \draw[FRMGold, -{Stealth[length=4.5pt]}, line width=0.9pt] (sw.east) -- (fl.west);
\end{tikzpicture}
\end{center}
\figcap{A pay-fixed, receive-floating swap is mapped as a short position in a coupon bond
alongside a long position in a floating-rate note.}

\subsection{Mapping approaches for nonlinear derivatives}
\begin{itemize}
\item \term{Option pricing.} While the option delta captures the option's sensitivity to the
      underlying price at a specific point, it is not constant. For options, the relationship is
      inherently non-linear.
\item \term{Short horizons.} The delta-normal method can be a reasonable approximation for short
      timeframes (e.g.\ one day) when the delta is relatively stable. For long-dated options,
      delta-normal is not accurate.
\end{itemize}

% ==========================================================
